\exercisesection{2}

\begin{enumerate}
\def\labelenumi{\arabic{enumi}.}
\item
  Let \(\Omega\) be an extension of a field K, E and F two extensions of K contained in \(\Omega\) and linearly disjoint over K and L an extension of K contained in E. Let \(f\) be a place of E with values in L which reduces to the identity automorphism on K. Show that there exists a unique place of K(E \(\cup\) F) with values in K(L \(\cup\) F), extending \(f\) and reducing to the identity automorphism on F. (If A is the ring of \(f\) , note that \(f\) can be extended uniquely to a homomorphism \(g\) from F{[}A{]} to K(L \(\cup\) F) reducing to the identity on F and that, if m is the ideal of \(f\) , every element \(z \in\) F{[}A{]} such that \(g(z) = 0\) can be written as \(xu\) , where \(x \in m\) , \(u \in\) F{[}A{]} and \(g(u) \neq 0\) . Conclude by noting that K(E \(\cup\) F) is the field of fractions of F{[}A{]}.)
\item
  Let K, K' be two extensions of a field k and f and \(f'\) places respectively of K and \(K'\) with values in the same algebraically closed field L. Suppose that the restrictions of f and \(f'\) to k coincide. Show that there exist a composite extension (F, i, i') of K and \(K'\) (Algebra, Chapter VIII, § 8) and a place g of F with values in L such that \(f(x) = g(i(x))\) for \(x \in \mathbf{K}\) and \(f'(x) = g(i'(x'))\) for \(x' \in \mathbf{K}'\) . (Let V, \(V'\) be the rings of f and \(f'\) , A their common intersection with k and h: \(V \otimes_{A} V' \to L\) the homomorphism such that \(h(a \otimes a') = f(a)f'(a')\) for \(a \in V\) , \(a' \in V'\) . Using the fact that V and \(V'\) are flat A-modules (§ 3, no. 5, Lemma 1), prove that, if \(a \neq 0\) , \(a' \neq 0\) , \(a \otimes a' = (a \otimes 1)(1 \otimes a')\) is not a divisor of zero; if S is the multiplicative subset of B = V \(\otimes_{A} V'\) consisting of these elements,
\end{enumerate}

\(S^{-1}B\) contains subfields \(K_{1}, K_{1}'\) respectively isomorphic to K and \(K'\) and is the ring generated by \(K_{1} \cup K_{1}'\) . Show that, if q is the kernel of h, there exists a prime ideal p in B such that \(p \subset q\) and \(p \cap S = \varnothing\) and show that F may be taken as the field of fractions of \(S^{-1}B/S^{-1}p\) .

\begin{enumerate}
\def\labelenumi{\arabic{enumi}.}
\setcounter{enumi}{2}
\tightlist
\item
  \begin{enumerate}
  \def\labelenumii{(\alph{enumii})}
  \tightlist
  \item
    Let \(\mathbf{K}\) be a field and \(f\) a place of \(\mathbf{K}\) with values in an orderable field \(\mathbf{L}\) (Algebra, Chapter VI, §2, Exercise 8); show that \(\mathbf{K}\) is orderable. (Note that, for a finite family of elements \((a_i)_{1\leq i\leq n}\) of \(\mathbf{K}\) , there is always an index \(j\) such that \(f(a_i / a_j)\neq \infty\) for all \(i\) .)
  \end{enumerate}
\end{enumerate}

If \((x_{i})_{1\leqslant i\leqslant n}\) is a sequence of elements of K such that \(f\left(\sum_{i = 1}^{n}x_{i}^{2}\right)\neq \infty\) , show that \(f(x_{i})\neq \infty\) for all \(i\)

\begin{enumerate}
\def\labelenumi{(\alph{enumi})}
\setcounter{enumi}{1}
\item
  Let K be an ordered field and G a subgroup of K; show that the ring \(F(G)\) of elements of K which are not infinitely great with respect to G (loc. cit., Exercise 11) is a valuation ring of K and the ideal \(I(G)\) of elements of \(F(G)\) infinitesimally small with respect to G is the maximal ideal of this ring; hence there is on K a place with values in \(k(G) = F(G)/I(G)\) , which reduces to the identity on G; this is called the canonical place of K with respect to G.
\item
  Show that, if there exists no extension of G contained in K, distinct from G and comparable with G, \(k(G)\) is algebraic over G (note that, if \(t \in K\) does not belong to G, the field \(G(t)\) contains an element \(u \neq 0\) infinitesimally small with respect to G and that \(G(t)\) is algebraic over \(G(u)\) ).
\end{enumerate}

¶ 4. An ordered field K is called Euclidean if, for all \(x \geqslant 0\) in K, there exists \(y \in \mathbf{K}\) such that \(x = y^2\) .

\begin{enumerate}
\def\labelenumi{(\alph{enumi})}
\item
  Let K be a Euclidean ordered field and G a subfield of K such that there exists no extension of G contained in K, distinct from G and comparable with G. Show that, if f is a place of K with values in an ordered field L, which is an algebraic extension of G, and reducing to the identity on G, then f is equivalent to the canonical place of K with respect to G. (f may be considered to take its values in a maximal ordered field N which is an algebraic extension of L. Observe that G is Euclidean and that, if \(x \in G\) , \(y \in K\) and x \textless{} y, then \(f(y) = \infty\) or \(f(y) > x\) . If A and p are the ring and ideal of the place f, show successively that \(p \subset I(G)\) , \(F(G) \subset A\) , \(I(G) \subset p\) and \(A \subset F(G)\) , noting that N is comparable with G.)
\item
  Let K be a Euclidean ordered field, f a place of K with values in a maximal ordered field L and A and p the ring and ideal of f.~Let G be a maximal subfield among the subfields E of K such that \(E \cap p = 0\) . Show that \(G \subset A\) and that G is algebraically closed in K and therefore Euclidean. Prove that there exists no extension of G contained in K, distinct from G and comparable with G. Moreover, the subfield of L generated by \(f(A)\) is algebraic over \(f(G)\) and f is therefore equivalent to the canonical place of K with respect to G.
\item
  Let K be a maximal ordered field, L a maximal ordered subfield of K, distinct from K and such that K is comparable with L (Algebra, Chapter VI, § 2, Exercise 15). Show that there exists no place f of K with values in L reducing to the identity automorphism on L.
\end{enumerate}

¶ 5. Let K be a field and f a place of K with values in a maximal ordered field L.

\begin{enumerate}
\def\labelenumi{(\alph{enumi})}
\item
  For all \(\alpha \in \mathbf{K}\) show that there exists an extension of \(f\) , either to \(\mathbf{K}(\sqrt{\alpha})\) or to \(\mathbf{K}(\sqrt{-\alpha})\) , which is a place with values in \(\mathbf{L}\) . (Consider two cases according to whether there exists \(a \in \mathbf{K}\) such that \(f(a^2\alpha)\) is neither 0 nor \(\infty\) or \(f(x^2\alpha)\) is equal to 0 or \(\infty\) for all \(x \in \mathbf{K}\) .)
\item
  Deduce from (a) that there exist an extension E of K which is a maximal ordered field and an extension of f to E which is a place of E with values in L. (Reduce it to the case where K is Euclidean and use Exercise 4 (a).)
\end{enumerate}

\begin{enumerate}
\def\labelenumi{\arabic{enumi}.}
\setcounter{enumi}{5}
\tightlist
\item
  Let A be an integrally closed domain, K its field of fractions, p a prime ideal of A and k the field of fractions of A/p.~Show that, if \(A_{p}\) is not a valuation ring, there exists a valuation ring V of K which dominates \(A_{p}\) and whose residue field is a transcendental extension of k (use Chapter V, §2, Exercise 14(b)).
\end{enumerate}

